\section{Discussão}
\label{sec:discussao}

\subsection{Contribuições Científicas}

A principal contribuição científica deste trabalho decorre da utilização de dados reais produzidos por uma planta automotiva em operação. Revisões da área de manutenção preditiva apontam que grande parte dos estudos é desenvolvida com dados simulados, bases laboratoriais ou conjuntos públicos já estruturados, enquanto aplicações documentadas com dados fabris reais permanecem menos frequentes~\cite{pdm4.0}. Restrições de confidencialidade, heterogeneidade dos sistemas industriais, ausência de padronização e dificuldade de validação operacional contribuem para essa lacuna. A pesquisa demonstra, portanto, como registros originados diretamente de um sistema industrial podem ser transformados em uma base analítica reproduzível, preservando limitações e descontinuidades que seriam menos evidentes em um conjunto previamente curado.

O estudo também se insere em um tema com literatura aplicada ainda restrita: os regimes de alarmes em cascata e seu impacto sobre a operação industrial. Trabalhos sobre manutenção preditiva frequentemente se concentram em sensores, classificação de falhas ou estimação de vida útil, enquanto há menos estudos dedicados à concentração temporal de alarmes e à antecipação de períodos nos quais recorrências rápidas dificultam a atuação dos operadores~\cite{rahaman2025alarmfloods}. A comparação direta com trabalhos semelhantes é ainda limitada pela escassez de bases públicas que preservem a ordem dos eventos, a identidade das máquinas e o contexto operacional necessário para caracterizar recorrências.

Diante da ausência de um rótulo externo de cascata, esta pesquisa propôs uma definição operacional e reprodutível baseada no tempo entre alarmes da mesma máquina. A contagem diária de eventos com IAT positivo inferior a 30 minutos não comprova propagação causal entre equipamentos, mas fornece uma medida objetiva da intensidade de recorrências rápidas. Essa definição permitiu converter um fenômeno operacional pouco estruturado em um problema supervisionado de regressão de contagem, tornando possível avaliar quantitativamente diferentes conjuntos de atributos e algoritmos.

Outra contribuição é a documentação integrada do percurso entre dado bruto e avaliação final. O trabalho compreende a identificação de hiatos, comparação entre séries bruta e sem \textit{outliers}, análise de recorrência por máquina e turno, avaliação espectral, decomposição, autocorrelação, seleção de variáveis exógenas, construção de atributos e validação temporal com teste final isolado. A explicitação dos dias com possível indisponibilidade e a manutenção de indicadores de qualidade impedem que todo zero seja interpretado automaticamente como operação normal. O código modular e os \textit{notebooks} executáveis tornam esse percurso auditável e adaptável a outros conjuntos de registros de alarmes.

\subsection{Implicações Práticas para a Indústria}

Embora o XGBoost completo tenha obtido o menor erro no teste final, o modelo simples apresenta uma relação particularmente favorável entre desempenho e custo de implantação. Sua versão com XGBoost alcançou MAE de 9,267, apenas 0,451 recorrência acima do modelo completo, utilizando 16 preditores em vez de 36. O conjunto simples depende exclusivamente do calendário, do histórico do volume global de alarmes e dos indicadores de \textit{gaps}; não exige atualização de séries macroeconômicas nem manutenção de agregações específicas para TR10, TR17 e TR19. Essa autonomia reduz dependências de integração e facilita sua incorporação à infraestrutura já utilizada para armazenar os alarmes da planta parceira.

Na prática, o \textit{pipeline} pode ser executado diariamente após o encerramento do dia operacional. A partir dos registros já coletados, calculam-se os volumes históricos e os indicadores de continuidade necessários para estimar a quantidade de recorrências rápidas esperada no dia seguinte. A previsão não identifica uma falha específica nem determina sua causa, mas pode funcionar como sinal antecipado de aumento da pressão operacional sobre as equipes responsáveis pelo acompanhamento dos alarmes.

Esse sinal pode apoiar o planejamento diário da manutenção preditiva. Dias com previsão elevada podem motivar reforço do monitoramento, revisão das ocorrências recentes, priorização de inspeções preventivas e adequação da disponibilidade de técnicos e recursos. A previsão também pode ser incorporada a um painel operacional, acompanhada dos indicadores de \textit{gaps}, para que o usuário saiba quando o histórico recente foi afetado por interrupções ou indisponibilidade de dados. Os limiares de alerta devem ser definidos em conjunto com especialistas da planta, considerando a capacidade da equipe e o custo relativo de falsos alertas e de picos não antecipados.

O modelo simples não substitui o diagnóstico de engenharia nem deve acionar intervenções automaticamente. Seu baixo $R^2$ e a subestimação dos maiores picos indicam que a previsão é mais adequada como informação complementar para priorização do que como estimativa precisa da carga máxima. Uma implementação piloto deve manter a decisão com os profissionais responsáveis, registrar previsões e eventos observados e acompanhar degradação de desempenho ao longo do tempo. A atualização periódica do modelo é necessária porque a análise temporal mostrou mudanças relevantes entre os regimes históricos da planta.

\subsection{Limitações e Ameaças à Validade}

A primeira limitação é a ausência de classificação de severidade dos alarmes. Todos os registros possuem a mesma representação na base, independentemente de exigirem intervenção imediata, serem resolvidos automaticamente ou não produzirem impacto mensurável na linha. Consequentemente, o alvo atribui o mesmo peso a recorrências operacionalmente distintas. Um dia com grande número de alarmes pouco críticos pode receber valor superior ao de um dia com poucas ocorrências severas, restringindo a interpretação da contagem como indicador direto de risco, custo ou necessidade de manutenção.

Também não existe um rótulo validado de regime de cascata. O limiar de IAT inferior a 30 minutos foi estabelecido a partir da distribuição observada e oferece uma definição consistente, mas permanece como \textit{proxy}. A base não permite demonstrar que um alarme causou outro, que ocorrências em máquinas diferentes pertencem ao mesmo processo de propagação ou que a concentração registrada produziu impacto sobre a produção. A validação dessa definição com especialistas da planta constitui condição importante para ampliar sua validade operacional.

A capacidade de generalização externa não foi testada. Os dados pertencem a uma única planta, a duas linhas da área \textit{Body-in-White} e a um intervalo histórico específico. O desempenho em outras plantas, processos produtivos, famílias de equipamentos ou períodos posteriores a maio de 2025 permanece desconhecido. Perfis de máquina, políticas de manutenção, calendários, turnos e práticas de registro podem alterar tanto a distribuição do alvo quanto a relevância das \textit{features}. Por isso, os resultados não devem ser transferidos diretamente para outro ambiente sem reconstrução do \textit{pipeline}, recalibração e nova validação temporal.

A validade temporal dentro da própria planta também é limitada pela não estacionariedade. As médias do alvo variaram intensamente entre os \textit{folds} expansivos, e o modelo vencedor apresentou $R^2$ negativo em todas as janelas de validação, apesar do resultado ligeiramente positivo no teste final. Os hiatos de origem parcialmente desconhecida reduzem a quantidade de observações elegíveis e podem coincidir com alterações operacionais não registradas. Embora os indicadores de \textit{gaps} preservem esse contexto, eles não recuperam a informação ausente nem determinam a causa da interrupção.

Por fim, a capacidade preditiva permanece limitada para eventos extremos. O XGBoost completo comprimiu as previsões em uma faixa muito inferior à amplitude observada e subestimou os dias de maior recorrência. A importância interna das árvores indica quais variáveis foram utilizadas nas divisões do modelo, mas não estabelece causalidade e pode distribuir relevância entre atributos correlacionados. A busca de hiperparâmetros foi deliberadamente compacta e a comparação ficou restrita a Random Forest e XGBoost, decisão compatível com o tempo disponível, mas que não esgota as abordagens possíveis para a série.

\section{Conclusões e Trabalhos Futuros}
\label{sec:conclusao}

\subsection{Conclusões}

O objetivo geral da pesquisa foi desenvolver e documentar uma metodologia de análise exploratória, engenharia de \textit{features} e modelagem preditiva aplicada a registros reais de alarmes industriais. Esse objetivo foi satisfeito. Os 88.411 eventos tratados foram organizados em uma série temporal contínua, investigados em suas dimensões temporal e operacional e convertidos em dois conjuntos de atributos para a previsão da quantidade diária de recorrências com IAT inferior a 30 minutos. A metodologia foi implementada em módulos reutilizáveis e exposta em \textit{notebooks}, permitindo reproduzir o percurso entre dados brutos, análise e avaliação final.

O cumprimento dos objetivos específicos pode ser sintetizado da seguinte forma:

\begin{enumerate}
    \item \textbf{Caracterização do conjunto de alarmes --- satisfeito.} A estrutura, o volume, o período de coleta, a distribuição entre linhas e máquinas e as principais limitações da base foram documentados. A análise confirmou a natureza discreta dos eventos, a ausência de severidade e a impossibilidade de inferir diretamente o impacto operacional de cada registro.
    \item \textbf{Diagnóstico e classificação dos hiatos --- parcialmente satisfeito.} Foram identificadas 48 sequências de dias com zero alarmes, das quais 32 possuem ao menos dois dias, e períodos conhecidos foram associados à pandemia, paradas de planta e recessos de calendário. Entretanto, os hiatos de 70 e 90 dias permaneceram sem causa validada, impedindo a classificação completa apenas com os dados disponíveis.
    \item \textbf{Análise de recorrência por máquina e turno --- satisfeito.} O IAT foi calculado dentro de cada máquina e permitiu identificar TR10 e TR17 pelo volume elevado e TR19 pela maior concentração de recorrências rápidas. As transições de turno e o intervalo entre o segundo e o terceiro turnos apresentaram comportamento temporal distinto.
    \item \textbf{Caracterização das sazonalidades e anomalias cíclicas --- satisfeito.} A FFT, a decomposição aditiva e as comparações por calendário foram realizadas nas séries bruta e tratada. Os resultados mostraram componentes temporais de múltiplas escalas, mas também evidenciaram que parte do padrão dominical e de outras aparentes sazonalidades é amplificada por episódios extremos e mudanças de regime.
    \item \textbf{Avaliação das variáveis exógenas --- satisfeito.} Selic e IPCA de bens duráveis apresentaram as associações mais consistentes e foram incorporados ao conjunto completo com disponibilidade mensal conservadora. PIB e balança comercial foram descartados por associações fracas ou instáveis.
    \item \textbf{Análise de autocorrelação --- satisfeito.} Foram examinadas ACF e PACF das séries bruta e sem \textit{outliers}, além da estabilidade anual dos coeficientes. A persistência de curto prazo e sinais semanais justificaram a escolha dos \textit{lags}, enquanto a divergência entre anos reforçou a necessidade de validação temporal.
    \item \textbf{Construção das \textit{features} --- satisfeito.} Foram obtidas matrizes simples e completa com 1.499 dias elegíveis e, respectivamente, 16 e 36 preditores. O desenho incorporou calendário, volumes defasados, histórico de \textit{gaps}, variáveis por máquina e, no conjunto completo, Selic e IPCA de duráveis, mantendo os dias sem registros sob tratamento explícito de disponibilidade.
    \item \textbf{Desenvolvimento e comparação dos modelos --- satisfeito com limitações de desempenho.} Random Forest e XGBoost foram avaliados nos dois conjuntos de atributos por validação expansiva e teste temporal isolado. O XGBoost completo obteve o melhor resultado, com MAE de 8,817, RMSE de 15,048 e $R^2$ de 0,051. A análise de importância destacou os volumes das máquinas selecionadas e o contexto dos \textit{gaps}, mas o baixo poder explicativo e a subestimação dos picos impedem interpretar o modelo como solução preditiva definitiva.
\end{enumerate}

Em síntese, o trabalho demonstrou que é possível extrair sinal preditivo de registros reais de alarmes utilizando apenas informações disponíveis antes do dia-alvo. O conjunto completo ofereceu o melhor desempenho estatístico, enquanto o conjunto simples preservou resultado próximo com menor custo de integração. A principal conclusão não é que os regimes de recorrência tenham se tornado plenamente previsíveis, mas que o histórico de alarmes, quando tratado com atenção às descontinuidades e à ordem temporal, fornece uma base mensurável para apoiar o monitoramento e orientar pesquisas posteriores.

\subsection{Proposta de Trabalhos Futuros}

Uma primeira extensão consiste em avaliar modelos de rede neural recorrente, especialmente LSTM~\cite{hochreiter1997lstm}, capazes de receber sequências temporais em vez de vetores de \textit{lags} previamente definidos. A comparação deverá utilizar os mesmos períodos de desenvolvimento, \textit{folds} expansivos e teste final, evitando que diferenças na divisão temporal favoreçam uma arquitetura. O uso de redes neurais também exigirá controle de complexidade, regularização e quantidade de dados suficiente para verificar se a maior flexibilidade produz ganho real sobre os modelos de árvores.

A generalização das \textit{features} e dos estimadores deve ser examinada com outros conjuntos de dados industriais semelhantes. Bases provenientes de outras linhas, plantas ou períodos posteriores permitiriam avaliar se os \textit{lags} globais, os indicadores de \textit{gaps} e a seleção de máquinas representam fenômenos transferíveis ou apenas características locais. A validação externa deverá contemplar recalibração do limiar de recorrência, comparação das distribuições e acompanhamento da degradação de desempenho.

Outra linha prioritária é aprofundar o estudo das variáveis nativas que não foram utilizadas nos modelos finais: \textit{fault codes}, posições, \textit{bilancellas} e tratos. Esses campos podem distinguir famílias de alarmes, regiões da linha e contextos operacionais associados a recorrências, mas sua inclusão depende de interpretação semântica, tratamento de alta cardinalidade e validação com especialistas da planta. Estudos de ablação devem medir o ganho incremental de cada grupo antes de ampliar o modelo implantado.

Também se recomenda construir uma classificação de severidade dos alarmes e obter rótulos operacionais associados a intervenções, paradas ou perda de produção. Com essas informações, o alvo poderá ponderar recorrências por criticidade e aproximar-se do impacto real sobre a manutenção. A definição de cascata baseada em IAT poderá então ser confrontada com avaliações de especialistas e com relações entre máquinas, permitindo distinguir mera concentração temporal de propagação operacional.

Por fim, modelos específicos para a assimetria e os eventos extremos devem ser investigados. Abordagens em duas etapas, com separação entre ocorrência de baixa atividade e previsão da intensidade positiva, funções de perda ponderadas ou modelos orientados a quantis podem reduzir a tendência de regressão à média observada neste trabalho. Paralelamente, um piloto prospectivo do modelo simples na planta parceira permitiria registrar previsões diariamente, definir limiares de alerta com a equipe de manutenção e avaliar se o sinal produzido contribui efetivamente para decisões operacionais.
